\textbf{Phần 1. Trắc nghiệm lựa chọn}
\setcounterpageref{ex}{0}

\begin{ex} %[4D1V02]
Tính giá trị của biểu thức sau bằng cách thuận tiện nhất:
$$28\,450 + 34\,700 + 21\,550 + 15\,300$$
\choice
{$90\,000$}
{\True $100\,000$}
{$105\,000$}
{$98\,000$}
\loigiai{
Ta nhóm các số hạng có tổng là số tròn chục nghìn:
\begin{align*}
& (28\,450 + 21\,550) + (34\,700 + 15\,300) \\
=& 50\,000 + 50\,000 = 100\,000
\end{align*}
\textbf{Đáp án của câu hỏi là B}
}
\end{ex}

\begin{ex} %[4D1V02]
Hiệu của hai số là $24\,500$. Nếu số bị trừ tăng thêm $3500$ đơn vị và số trừ giảm đi $1500$ đơn vị thì hiệu mới của hai số là
\choice
{$26\,500$}
{$22\,500$}
{$28\,000$}
{\True $29\,500$}
\loigiai{
Khi số bị trừ tăng thêm $3500$ đơn vị thì hiệu tăng thêm $3500$ đơn vị.\\
Khi số trừ giảm đi $1500$ đơn vị thì hiệu tăng thêm tiếp $1500$ đơn vị.\\
Hiệu mới của hai số là:
$$24\,500 + 3500 + 1500 = 29\,500$$
\textbf{Đáp án của câu hỏi là D}
}
\end{ex}

\begin{ex} %[4D1V02]
Một số khi nhân với $6$ rồi trừ đi $18\,400$ thì được kết quả là $44\,600$. Số đó là
\choice
{\True $10\,500$}
{$11\,500$}
{$10\,200$}
{$9500$}
\loigiai{
Gọi số cần tìm là $x$. Theo bài ra ta có:
\begin{align*}
x \times 6 - 18\,400 &= 44\,600 \\
x \times 6 &= 44\,600 + 18\,400 \\
x \times 6 &= 63\,000 \\
x &= 63\,000 : 6 \\
x &= 10\,500
\end{align*}
\textbf{Đáp án của câu hỏi là A}
}
\end{ex}

\begin{ex} %[4D1V02]
Trong một phép chia có số chia là $8$, thương là $11\,425$ và số dư là số dư lớn nhất có thể. Số bị chia của phép chia đó là
\choice
{$91\,400$}
{$91\,408$}
{\True $91\,407$}
{$91\,405$}
\loigiai{
Vì số chia là $8$ nên số dư lớn nhất có thể có là $8 - 1 = 7$.\\
Số bị chia là:
$$11\,425 \times 8 + 7 = 91\,400 + 7 = 91\,407$$
\textbf{Đáp án của câu hỏi là C}
}
\end{ex}

\begin{ex} %[4D1V02]
Giá trị của biểu thức $80\,000 - 60\,000 : (2 \times 3)$ là
\choice
{$10\,000$}
{\True $70\,000$}
{$30\,000$}
{$50\,000$}
\loigiai{
Thực hiện phép tính trong ngoặc trước, sau đó đến phép chia và phép trừ:
\begin{align*}
80\,000 - 60\,000 : (2 \times 3) &= 80\,000 - 60\,000 : 6 \\
&= 80\,000 - 10\,000 \\
&= 70\,000
\end{align*}
\textbf{Đáp án của câu hỏi là B}
}
\end{ex}

\begin{ex} %[4D1V02]
Tìm số tự nhiên $x$, biết: $(x + 15\,200) : 4 = 12\,300$.
\choice
{\True $34\,000$}
{$33\,000$}
{$35\,000$}
{$49\,200$}
\loigiai{
Ta có:
\begin{align*}
(x + 15\,200) : 4 &= 12\,300 \\
x + 15\,200 &= 12\,300 \times 4 \\
x + 15\,200 &= 49\,200 \\
x &= 49\,200 - 15\,200 \\
x &= 34\,000
\end{align*}
\textbf{Đáp án của câu hỏi là A}
}
\end{ex}

\begin{ex} %[4D1V02]
Một siêu thị nhập về ba chuyến xe chở gạo. Chuyến thứ nhất chở $15\,800\text{ kg}$ gạo. Chuyến thứ hai chở số gạo gấp đôi chuyến thứ nhất. Chuyến thứ ba chở ít hơn chuyến thứ hai $4600\text{ kg}$ gạo. Cả ba chuyến xe chở được tất cả số ki-lô-gam gạo là
\choice
{$68\,400\text{ kg}$}
{$72\,200\text{ kg}$}
{$75\,000\text{ kg}$}
{\True $74\,400\text{ kg}$}
\loigiai{
Chuyến thứ hai chở số gạo là:
$$15\,800 \times 2 = 31\,600\text{ (kg)}$$
Chuyến thứ ba chở số gạo là:
$$31\,600 - 4600 = 27\,000\text{ (kg)}$$
Cả ba chuyến xe chở được tất cả số gạo là:
$$15\,800 + 31\,600 + 27\,000 = 74\,400\text{ (kg)}$$
\textbf{Đáp án của câu hỏi là D}
}
\end{ex}

\begin{ex} %[4D1V02]
Tính giá trị của biểu thức: $A = 8500 \times 4 + 8500 \times 6$.
\choice
{$80\,000$}
{$82\,500$}
{\True $85\,000$}
{$90\,000$}
\loigiai{
Áp dụng tính chất nhân một số với một tổng:
\begin{align*}
A &= 8500 \times 4 + 8500 \times 6 \\
&= 8500 \times (4 + 6) \\
&= 8500 \times 10 \\
&= 85\,000
\end{align*}
\textbf{Đáp án của câu hỏi là C}
}
\end{ex}

\begin{ex} %[4D1V02]
Để hoàn thành một đoạn kênh mương trong $6$ ngày thì cần $15$ công nhân. Nếu muốn hoàn thành đoạn kênh mương đó sớm hơn $1$ ngày với cùng năng suất làm việc của mỗi người thì cần bổ sung thêm bao nhiêu công nhân?
\choice
{$18$ công nhân}
{\True $3$ công nhân}
{$5$ công nhân}
{$6$ công nhân}
\loigiai{
Để hoàn thành công việc trong $1$ ngày cần số công nhân là:
$$15 \times 6 = 90\text{ (công nhân)}$$
Để hoàn thành công việc sớm hơn $1$ ngày (tức trong $6 - 1 = 5$ ngày) cần số công nhân là:
$$90 : 5 = 18\text{ (công nhân)}$$
Số công nhân cần bổ sung thêm là:
$$18 - 15 = 3\text{ (công nhân)}$$
\textbf{Đáp án của câu hỏi là B}
}
\end{ex}

\begin{ex} %[4D1V02]
Một phép chia cho $9$ có số dư là $6$. Hỏi cần thêm vào số bị chia ít nhất bao nhiêu đơn vị để phép chia trở thành phép chia hết và thương tăng thêm $2$ đơn vị?
\choice
{$3$ đơn vị}
{$9$ đơn vị}
{$15$ đơn vị}
{\True $12$ đơn vị}
\loigiai{
Số dư ban đầu là $6$. Để phép chia chia hết và thương tăng thêm $1$ đơn vị, cần thêm vào số bị chia ít nhất:
$$9 - 6 = 3\text{ (đơn vị)}$$
Để thương tăng thêm $2$ đơn vị thì cần thêm tiếp một lần số chia nữa, tức là thêm:
$$3 + 9 = 12\text{ (đơn vị)}$$
\textbf{Đáp án của câu hỏi là D}
}
\end{ex}

\begin{ex} %[4D1V02]
Bác An mua $3\text{ kg}$ cam và $2\text{ kg}$ táo hết $95\,000$ đồng. Bác Bình mua $3\text{ kg}$ cam và $5\text{ kg}$ táo cùng loại hết $155\,000$ đồng. Giá tiền của $1\text{ kg}$ táo là
\choice
{\True $20\,000$ đồng}
{$15\,000$ đồng}
{$25\,000$ đồng}
{$18\,000$ đồng}
\loigiai{
Bác Bình mua nhiều hơn bác An số ki-lô-gam táo là:
$$5 - 2 = 3\text{ (kg)}$$
Số tiền mua $3\text{ kg}$ táo đó là:
$$155\,000 - 95\,000 = 60\,000\text{ (đồng)}$$
Giá tiền của $1\text{ kg}$ táo là:
$$60\,000 : 3 = 20\,000\text{ (đồng)}$$
\textbf{Đáp án của câu hỏi là A}
}
\end{ex}

\begin{ex} %[4D1V02]
Tìm số tự nhiên bé nhất có 5 chữ số chia hết cho cả $2$, $5$ và $9$.
\choice
{$10\,800$}
{$10\,008$}
{\True $10\,080$}
{$11\,700$}
\loigiai{
Một số chia hết cho cả 2 và 5 thì chữ số hàng đơn vị phải là $0$.\\
Số cần tìm có dạng $\overline{abcd0}$ với $a \neq 0$.\\
Để số đó là bé nhất có 5 chữ số thì ta chọn các chữ số từ trái sang phải nhỏ nhất có thể:\\
- Chọn $a = 1, b = 0, c = 0$.\\
Khi đó số có dạng $\overline{100d0}$. Tổng các chữ số là: $1 + 0 + 0 + d + 0 = 1 + d$.\\
Để số chia hết cho 9 thì $1 + d$ phải chia hết cho 9, do đó $d = 8$.\\
Vậy số tự nhiên bé nhất thỏa mãn là $10\,080$.\\
\textbf{Đáp án của câu hỏi là C}
}
\end{ex}

\textbf{Phần 2. Trắc nghiệm đúng sai}
\setcounterpageref{ex}{0}

\begin{ex} %[4D1V02]
Xét tính đúng sai của các khẳng định sau về tính chất của các phép tính:
\choiceTF
{\True Trong phép chia có dư với số chia bằng $7$, số dư lớn nhất có thể có là $6$.}
{Nếu mỗi số hạng trong một tổng cùng tăng thêm $2500$ đơn vị thì tổng đó sẽ tăng thêm $2500$ đơn vị.}
{\True Tích của một số bất kì với số $0$ luôn bằng $0$.}
{\True Khi chia số $64\,048$ cho $8$ ta được kết quả là $8006$.}
\loigiai{
\begin{itemize}
    \item Số dư luôn luôn nhỏ hơn số chia. Vì số chia là 7 nên số dư lớn nhất có thể là $7 - 1 = 6$. Khẳng định này Đúng.
    \item Nếu mỗi số hạng trong tổng cùng tăng thêm $2500$ đơn vị thì tổng sẽ tăng thêm một lượng bằng $2500$ nhân với số lượng các số hạng. Do đó khẳng định này Sai.
    \item Theo tính chất nhân với số 0, mọi số nhân với 0 đều bằng 0. Khẳng định này Đúng.
    \item Thực hiện phép chia: $64\,048 : 8 = 8006$. Khẳng định này Đúng.
\end{itemize}
\begin{itemize}
    \item a) \textbf{Đúng}.
    \item b) \textbf{Sai}.
    \item c) \textbf{Đúng}.
    \item d) \textbf{Đúng}.
\end{itemize}
}
\end{ex}

\begin{ex} %[4D1V02]
Một cửa hàng vật liệu có ba kho chứa xi măng. Kho A chứa $24\,500\text{ kg}$ xi măng. Kho B chứa nhiều hơn kho A là $5000\text{ kg}$ xi măng. Kho C chứa số xi măng bằng một nửa tổng số xi măng của cả hai kho A và B. Xét tính đúng sai của các khẳng định sau:
\choiceTF
{\True Kho B chứa $29\,500\text{ kg}$ xi măng.}
{Kho C chứa $26\,000\text{ kg}$ xi măng.}
{\True Tổng số xi măng của cả ba kho là $81\,000\text{ kg}$.}
{\True Kho C chứa nhiều hơn kho A là $2500\text{ kg}$ xi măng.}
\loigiai{
\begin{itemize}
    \item Số xi măng ở kho B là: $24\,500 + 5000 = 29\,500\text{ (kg)}$. Khẳng định a Đúng.
    \item Tổng số xi măng của kho A và kho B là: $24\,500 + 29\,500 = 54\,000\text{ (kg)}$.\\
    Số xi măng ở kho C là: $54\,000 : 2 = 27\,000\text{ (kg)}$ (không phải $26\,000\text{ kg}$). Khẳng định b Sai.
    \item Tổng số xi măng của cả ba kho là: $54\,000 + 27\,000 = 81\,000\text{ (kg)}$. Khẳng định c Đúng.
    \item Kho C nhiều hơn kho A số xi măng là: $27\,000 - 24\,500 = 2500\text{ (kg)}$. Khẳng định d Đúng.
\end{itemize}
\begin{itemize}
    \item a) \textbf{Đúng}.
    \item b) \textbf{Sai}.
    \item c) \textbf{Đúng}.
    \item d) \textbf{Đúng}.
\end{itemize}
}
\end{ex}

\textbf{Phần 3. Trả lời ngắn}
\setcounterpageref{ex}{0}

\begin{ex} %[4D1V02]
Tính giá trị của biểu thức bằng cách thuận tiện nhất:
$$14\,500 + 27\,800 + 35\,500 + 12\,200$$
\shortans{90000}
\loigiai{
\textbf{Đáp số:} $90000$\\
Nhóm các số hạng tạo thành số tròn chục nghìn:
\begin{align*}
& (14\,500 + 35\,500) + (27\,800 + 12\,200) \\
=& 50\,000 + 40\,000 = 90\,000
\end{align*}
}
\end{ex}

\begin{ex} %[4D1V02]
Tìm số bị chia trong một phép chia có số chia là $7$, thương là số lẻ bé nhất có 5 chữ số khác nhau và số dư là số dư lớn nhất có thể có.
\shortans{71651}
\loigiai{
\textbf{Đáp số:} $71651$\\
- Số lẻ bé nhất có 5 chữ số khác nhau là $10\,235$.\\
- Vì số chia là $7$ nên số dư lớn nhất có thể có là $7 - 1 = 6$.\\
Số bị chia cần tìm là:
$$10\,235 \times 7 + 6 = 71\,645 + 6 = 71\,651$$
}
\end{ex}

\begin{ex} %[4D1V02]
Tìm giá trị của $x$, biết: $(x - 14\,500) \times 3 = 45\,000$.
\shortans{29500}
\loigiai{
\textbf{Đáp số:} $29500$\\
Ta có:
\begin{align*}
(x - 14\,500) \times 3 &= 45\,000 \\
x - 14\,500 &= 45\,000 : 3 \\
x - 14\,500 &= 15\,000 \\
x &= 15\,000 + 14\,500 \\
x &= 29\,500
\end{align*}
}
\end{ex}

\begin{ex} %[4D1V02]
Hai thùng chứa tất cả $36\,800\text{ l}$ dầu. Nếu chuyển $2400\text{ l}$ dầu từ thùng thứ nhất sang thùng thứ hai thì số dầu ở hai thùng bằng nhau. Hỏi ban đầu thùng thứ nhất chứa bao nhiêu lít dầu?
\shortans{20800}
\loigiai{
\textbf{Đáp số:} $20800$\\
Khi chuyển $2400\text{ l}$ dầu từ thùng thứ nhất sang thùng thứ hai thì hai thùng bằng nhau, vậy thùng thứ nhất nhiều hơn thùng thứ hai số lít dầu là:
$$2400 \times 2 = 4800\text{ (l)}$$
Ban đầu thùng thứ nhất chứa số lít dầu là:
$$(36\,800 + 4800) : 2 = 20\,800\text{ (l)}$$
}
\end{ex}

\begin{ex} %[4D1V02]
Một xưởng may trong $4$ ngày may được $1800$ chiếc áo. Hỏi với năng suất may như vậy, xưởng đó cần bao nhiêu ngày để may được $3150$ chiếc áo?
\shortans{7}
\loigiai{
\textbf{Đáp số:} $7$\\
Mỗi ngày xưởng đó may được số chiếc áo là:
$$1800 : 4 = 450\text{ (chiếc)}$$
Để may được $3150$ chiếc áo thì xưởng đó cần số ngày là:
$$3150 : 450 = 7\text{ (ngày)}$$
}
\end{ex}

\begin{ex} %[4D1V02]
Tính tổng của dãy số cách đều sau:
$$10\,100 + 10\,200 + 10\,300 + \dots + 10\,900$$
\shortans{94500}
\loigiai{
\textbf{Đáp số:} $94500$\\
Số lượng các số hạng của dãy số là:
$$(10\,900 - 10\,100) : 100 + 1 = 9\text{ (số)}$$
Tổng của dãy số là:
$$(10\,100 + 10\,900) \times 9 : 2 = 21\,000 \times 9 : 2 = 94\,500$$
}
\end{ex}
